\question[4]
Erkläre in ein bis zwei Sätzen.
\begin{parts}
\part Was ist der Unterschied zwischen einer Klasse und einem Objekt?
\part Welche Aufgabe hat der Konstruktor?
\part Was bedeutet Geheimnisprinzip?
\part Was ist der Unterschied zwischen einer Assoziation und einer Vererbung?
\end{parts}
\begin{solutionbox}{3.5cm}
a) Eine Klasse ist der Bauplan, ein Objekt ist eine nach diesem Bauplan erzeugte Instanz. \\
b) Er wird beim Erzeugen eines Objekts durchlaufen und versetzt es in einen vernünftigen Anfangszustand. \\
c) Die Attribute sind nach außen nicht direkt sichtbar, der Zugriff läuft über Methoden. Wie die Klasse ihre
Daten speichert, bleibt verborgen. \\
d) Assoziation: „hat ein“, ein Attribut verweist auf ein Objekt einer anderen Klasse. Vererbung: „ist ein“, die
Unterklasse übernimmt die Attribute und Methoden der Oberklasse.
\end{solutionbox}
